%=====================================================================
%  DEMOSTRACIONES DE CAMPOS CONOCIDOS
%  Distribuciones de cargas continuas
%  CFIS 344 - CFSA 3440 - CFSV 344 - Electricidad y Magnetismo
%  Universidad Andres Bello
%  Apoyo docente: Javier Patino
%
%  Compilar:  pdflatex Demostraciones_Campos_Conocidos   (dos veces)
%=====================================================================
\documentclass[a4paper,justified,nobib,openany]{tufte-book}

%---------------------------------------------------------------------
% 1. Idioma y codificacion
%---------------------------------------------------------------------
\usepackage[utf8]{inputenc}
\usepackage[T1]{fontenc}
\usepackage[spanish,es-noquoting,es-nolists,es-tabla]{babel}

%---------------------------------------------------------------------
% 2. Matematicas
%---------------------------------------------------------------------
\usepackage{amsmath,amssymb}
\usepackage{bm}
\usepackage{cancel}

%---------------------------------------------------------------------
% 3. Graficos, color y tablas
%---------------------------------------------------------------------
\usepackage{graphicx}
\graphicspath{{figuras/}{./}}
\usepackage[table]{xcolor}
\usepackage{booktabs}
\usepackage{array}
\usepackage{tikz}
\usetikzlibrary{arrows.meta,positioning,shapes.geometric,calc,
                decorations.markings,decorations.pathreplacing,angles,quotes}

% --- Paleta institucional UNAB ---
\definecolor{unabazul}{HTML}{06304B}
\definecolor{unabrojo}{HTML}{9A1B34}
\definecolor{unabgris}{HTML}{4A4A4A}
\definecolor{cajafondo}{HTML}{F2F6F9}
\definecolor{ambarfondo}{HTML}{FFF7E6}
\definecolor{ambarborde}{HTML}{B8860B}
\definecolor{verdefondo}{HTML}{F0F6F0}
\definecolor{verdeborde}{HTML}{3F6B46}
\definecolor{cargapos}{HTML}{C0392B}
\definecolor{carganeg}{HTML}{1F4E79}
\definecolor{lineacampo}{HTML}{3B6FA0}
\definecolor{acento}{HTML}{1F6FB2}

%---------------------------------------------------------------------
% 4. Minicuadros
%---------------------------------------------------------------------
\usepackage[most]{tcolorbox}
\tcbuselibrary{skins,breakable}

\tcbset{
  cuadrobase/.style={
    enhanced, breakable, sharp corners,
    boxrule=0pt, leftrule=2.5pt,
    left=8pt, right=8pt, top=6pt, bottom=6pt,
    fonttitle=\sffamily\bfseries\footnotesize,
    coltitle=white,
    attach boxed title to top left={xshift=-2.5pt, yshift=-2pt},
    boxed title style={sharp corners, boxrule=0pt, left=6pt, right=6pt,
                       top=2.5pt, bottom=2.5pt},
    before skip=11pt, after skip=11pt
  }
}

\newtcolorbox{ideaclave}[1][Idea clave]{%
  cuadrobase, colback=cajafondo, colframe=unabazul,
  colbacktitle=unabazul, title={#1}}

\newtcolorbox{historia}[1][Nota histórica]{%
  cuadrobase, colback=white, colframe=unabrojo,
  colbacktitle=unabrojo, title={#1}}

\newtcolorbox{ojo}[1][¡Ojo!]{%
  cuadrobase, colback=ambarfondo, colframe=ambarborde,
  colbacktitle=ambarborde, title={#1}}

\newtcolorbox{dato}[1][Comprobación]{%
  cuadrobase, colback=verdefondo, colframe=verdeborde,
  colbacktitle=verdeborde, title={#1}}

\newtcolorbox{herramienta}[1][Herramienta]{%
  cuadrobase, colback=white, colframe=acento,
  colbacktitle=acento, title={#1}}

%---------------------------------------------------------------------
% 5. Enlaces
%---------------------------------------------------------------------
\usepackage{hyperref}
\hypersetup{
  colorlinks = true,
  linkcolor  = unabazul,
  citecolor  = unabrojo,
  urlcolor   = unabrojo,
  linktoc    = all,
  bookmarksnumbered = true,
  bookmarksopen     = true,
  bookmarksopenlevel= 1,
  pdftitle   = {Demostraciones de campos conocidos. Distribuciones continuas},
  pdfauthor  = {Javier Patiño},
  pdfsubject = {Electricidad y Magnetismo, Universidad Andrés Bello}
}

%---------------------------------------------------------------------
% 6. Estilo de titulos
%---------------------------------------------------------------------
\addto\captionsspanish{\renewcommand{\chaptername}{Demostración}}
\setcounter{secnumdepth}{2}
\renewcommand{\thesection}{\arabic{section}}
\renewcommand{\thesubsection}{\arabic{section}.\arabic{subsection}}

\titleformat{\chapter}%
  [display]%
  {\relax\normalfont\LARGE\itshape\color{unabazul}}%
  {}%
  {0pt}%
  {}%
  [{\vspace{3pt}\color{unabrojo}\titlerule[1.1pt]}]
\titlespacing*{\chapter}{0pt}{0pt}{20pt}

\titleformat{\section}%
  [hang]%
  {\normalfont\Large\itshape\color{unabazul}}%
  {\thesection.}%
  {0.55em}%
  {}
\titlespacing*{\section}{0pt}{1.9\baselineskip}{0.55\baselineskip}

\titleformat{\subsection}%
  [hang]%
  {\normalfont\large\itshape\color{unabgris}}%
  {\thesubsection.}%
  {0.55em}%
  {}
\titlespacing*{\subsection}{0pt}{1.3\baselineskip}{0.35\baselineskip}

%---------------------------------------------------------------------
% 7. Atajos propios
%---------------------------------------------------------------------
\newcommand{\nota}[1]{\sidenote{#1}}
\newcommand{\clave}[1]{\textcolor{unabrojo}{\textbf{#1}}}
\newcommand{\resultado}[1]{%
  \begin{center}%
  \setlength{\fboxsep}{7pt}%
  \fcolorbox{unabrojo}{ambarfondo}{$\displaystyle #1$}%
  \end{center}}

% Vectores y versores
\renewcommand{\vec}[1]{\mathbf{#1}}
\newcommand{\vE}{\mathbf{E}}
\newcommand{\vF}{\mathbf{F}}
\newcommand{\vr}{\mathbf{r}}
\newcommand{\versor}[1]{\hat{\mathbf{#1}}}
\newcommand{\versorg}[1]{\hat{\bm{#1}}}
\newcommand{\norm}[1]{\left\lVert #1 \right\rVert}

% Etiqueta de paso dentro de una demostracion
\newcommand{\paso}[1]{%
  \medskip\noindent
  \textcolor{unabrojo}{\sffamily\bfseries\small Paso #1.}\ }

\usepackage{enumitem}
\setlist{itemsep=2pt, topsep=4pt, parsep=0pt}

%=====================================================================
\begin{document}
%=====================================================================

%---------------------------------------------------------------------
% PORTADA
%---------------------------------------------------------------------
\begin{titlepage}
\thispagestyle{empty}
\begin{tikzpicture}[remember picture, overlay]
  \fill[unabazul] (current page.north west) rectangle
        ([xshift=2.1cm]current page.south west);
  \fill[unabrojo] ([xshift=2.1cm]current page.north west) rectangle
        ([xshift=2.35cm]current page.south west);
  \node[anchor=north east, inner sep=0pt]
        at ([xshift=-2.0cm, yshift=-2.0cm]current page.north east)
        {\includegraphics[width=3.4cm]{unab.png}};
\end{tikzpicture}

\vspace*{5.0cm}
\hspace*{1.0cm}
\begin{minipage}{0.80\textwidth}
  \raggedright
  {\sffamily\footnotesize\bfseries\color{unabrojo}
   UNIVERSIDAD ANDRÉS BELLO\par}
  {\sffamily\footnotesize\color{unabgris}
   FACULTAD DE CIENCIAS EXACTAS\par
   DEPARTAMENTO DE FÍSICA Y ASTRONOMÍA\par}

  \vspace{1.2cm}
  {\color{unabazul}\Huge\itshape Demostraciones de\\[6pt]
   campos conocidos\par}

  \vspace{0.5cm}
  {\color{unabgris}\Large\itshape Distribuciones de cargas continuas\par}

  \vspace{0.6cm}
  {\color{unabrojo}\rule{10.5cm}{1.2pt}\par}
  \vspace{0.45cm}

  {\sffamily\normalsize\color{unabgris}
   CFIS 344 · CFSA 3440 · CFSV 344 — Electricidad y Magnetismo\par}

  \vspace{1.8cm}
  {\large\itshape Javier Patiño\par}
  \vspace{2pt}
  {\footnotesize\color{unabgris} Apoyo docente\par}
  \vspace{6pt}
  {\footnotesize\color{unabgris}
   Profesor responsable: Rodrigo Vergara Rojas\par
   Universidad Andrés Bello, Santiago de Chile\par}

  \vspace{0.9cm}
  {\footnotesize\color{unabgris}
   Material de apoyo para las grabaciones de la Semana 2.\par
   \today\par}
\end{minipage}

\begin{tikzpicture}[remember picture, overlay]
  % anillo decorativo con su campo axial, al pie
  \begin{scope}[shift={([xshift=-4.6cm, yshift=3.6cm]current page.south east)},
                scale=0.95]
    \draw[lineacampo, thick] (0,0) ellipse (1.6 and 0.50);
    \foreach \a in {0,45,...,315}{
      \filldraw[cargapos] ({1.6*cos(\a)},{0.50*sin(\a)}) circle (0.075);}
    \draw[unabgris!60, dashed] (0,0) -- (0,2.0);
    \draw[-{Latex[length=6pt]}, unabrojo, very thick] (0,1.55) -- (0,2.45);
    \filldraw[unabazul] (0,1.55) circle (0.055);
    \node[right, font=\scriptsize, unabgris] at (0.12,1.55) {$P$};
  \end{scope}
\end{tikzpicture}
\end{titlepage}

%---------------------------------------------------------------------
% INDICE
%---------------------------------------------------------------------
\setcounter{tocdepth}{1}
\tableofcontents
\clearpage

%=====================================================================
%  CUERPO
%=====================================================================
\input{secciones/00_metodo}
\input{secciones/01_arco}
\input{secciones/02_anillo}
\input{secciones/03_disco}
\input{secciones/04_plano}
\input{secciones/05_barra_eje}
\input{secciones/06_barra_mediatriz}
\input{secciones/07_resumen}

\end{document}
